
\section{Appendix A: Extremal length and width}\label{apppendix: width}
[If you have come to this section from  \cite{decorations} or \cite{molecules},
looking for Theorem 9.1 of this paper,
it is now known as Theorem \ref{improving of length}.]

There is a wealth of sources containing background material on extremal length,
see, e.g., the book of Ahlfors \cite{A}.
We will briefly summarize the necessary minimum
(see also the Appendix of \cite{covering lemma}).

\subsection{Definitions} \label{A:def}
In this section we will think of a path as a smooth map $\gamma\from I \to U$,
where $I \subset R$ is an interval (that may be open or closed),
and $\gamma$ is not necessarily proper.

Let $\GG$ be a family of paths on a Riemann surface $U$.
Given a (measurable) conformal metric $\mu= \mu(z) |dz|$ on $U$,
let
$$   \mu (\GG) = \inf_{\gamma\in \GG}  \mu (\gamma).$$
 where $\mu(\gamma) $ stands for the $\mu$-length of $\gamma$.
The length of $\GG$  with respect to $\mu$ is defined as
$$
   \LL_\mu (\GG) = \frac {\mu(\GG)^2 } {\area_\mu (U)},
$$
where $\area_\mu$ is an area form of $\mu$.
Taking the supremum over all conformal metrics $\mu$, we obtain the {\it extremal length}
$\LL(\GG)$ of the family $\GG$.

  The {\it extremal width} is the inverse of the extremal length:
$$
 \WW(\GG)=\LL^{-1}(\GG).
$$
It can be also defined as follows. Consider all conformal  metrics $\mu$ such that
$\mu(\gamma)\geq 1$ for any $\gamma\in \GG$.  Then $\WW(\GG)$ is the infimum of
the areas $\area_\mu(U)$ of  all such metrics.


\subsection{Electric circuit laws}\label{electric}
We say that a family  $\GG$ of paths {\it overflows} a family $\GG'$ if
any path of $\GG$ contains some path of $\GG$.
Let us say that $\GG$ {\it disjointly overflows}  two families, $\GG_1$ and $\GG_2$,
if any path of $\gamma\in \GG$ contains the disjoint union $\gamma_1\sqcup \gamma_2$
of two paths $\gamma_i\in \HH_i$.

The following crucial properties of the extremal length and width show
that the former behaves like the resistance in electric circuits,
while the latter behaves like conductance.

\begin{serieslaw}
Let $\GG$ be a family of paths that disjointly overflows two other families,
                      $\GG_1$ and $\GG_2$.
Then
$$  \LL(\GG) \geq \LL(\GG_1) + \LL(\GG_2), $$
or equivalently,
$$  \WW(\GG)\leq \WW(\GG_1)\bigoplus \WW(\GG_2). $$
\end{serieslaw}

\begin{parallellaw}
  For any two families $\GG_1$ and $\GG_2$ of paths we have:
$$
    \WW(\GG_1\cup \GG_2)\leq \WW(\GG_1) + \WW(\GG_2).
$$
If $\GG_1$ and $\GG_2$ are contained in two disjoint open sets, then
$$
  \WW(\GG_1\cup \GG_2) = \WW(\GG_1) + \WW(\GG_2)
$$
\end{parallellaw}



\subsection{Transformation rules}\label{trans rules sec}
 Both extremal length and extremal width  are conformal invariants.
 More generally, we have:
\begin{lem}\label{increase}
   Let $f\from U\to V$ be a holomorphic map between two Riemann surfaces,
and let $\GG$ be a family of paths on $U$. Then
$$
      \LL(f(\GG))\geq \LL(\GG).
$$
\end{lem}
See Lemma 4.1 of \cite{covering lemma} for a proof.
\begin{rem}
We push forward \emph{parametrized} paths, so each $\gamma \in \GG$ is mapped to $f \circ \gamma \in f(\GG)$. 
For example, if $f\from U\to V$ is a degree 2 covering map between annuli (of finite modulus), and $\GG$ is the family of closed paths that go around $U$,
then $f(\GG)$ includes the closed paths that go twice around another closed path that goes once around $V$. 
While the extremal length of the closed paths that go once around $V$ is $\LL(\GG)/2$,
we actually have $\LL(f(\GG)) = 2 \LL(\GG)$ in this example.
\end{rem}

\begin{cor}\label{two families}
     Under the circumstances of the previous lemma,
let $\HH$ be a family of paths in $V$ satisfying the following lifting
property: any path $\gamma\in \HH$ contains an arc that lifts to some path
in $\GG$. Then $\LL(\HH)\geq \LL(\GG)$.
\end{cor}

\begin{proof}
   The lifting property means that the family $\HH$
overflows the family $f(\GG)$. Hence $\LL(\HH)\geq \LL(f(\GG))$,
and the conclusion follows.
\end{proof}

\begin{cor}\label{Q}
  Let $Q$ and $Q'$ be two quadrilaterals,
and let $e\from Q\to Q'$ be a holomorphic map that maps the horizontal sides of $Q$ to the horizontal sides of $Q'$.
Then $\WW(Q)\leq \WW(Q')$.
\end{cor}

\begin{proof}
  Let $\GG$ (resp. $\GG'$) be the family of horizontal paths in $Q$ (resp., in $Q'$).
 Since the horizontal sides of $Q$ are mapped to the horizontal sides of $Q'$,
these families satisfy the lifting property of the previous Corollary.
Hence $\LL(\GG)\leq \LL(\GG')$, and we are done.
\end{proof}

\begin{lem}\label{modulus transform-2}
 Let $f\from U\to V$ be a branched covering between two compact Riemann surfaces with boundary.
Let $A$ be an archipelago in $U$,  $B=f(A)$, and
assume that  $f\from  A\to B$ is a branched covering of degree $d$. Then
$$
      \mod(V, B) \geq d\, \mod (U,  A).
$$
\end{lem}

See Lemma 4.3 of \cite{covering lemma} for a proof.


Given two compact subsets $A$ and $B$ in a Riemann surface $S$,
let $\WW_S (A,B)$ stand for the {\it extremal width}
between them, i.e., the extremal width of the family of paths connecting $A$ to $B$.

\begin{lem}\label{modulus transform-1}
Let $S$ and $S'$ be two compact Riemann surfaces with boundary.
and let  $f\from S\to S'$ be a branched covering of degree $D$.
Let $S'=A'\sqcup B'$, where $A'$ and $B'$ are closed,
and let  $A=f^{-1}(A')$, $B=f^{-1}(B')$.
Then
$$
      \WW_S(A, B)= D \, \WW (A', B').
$$
\end{lem}

See \cite{A} for a proof.
It makes use of the fact that
the extremal width $\WW(A,B)$ is achieved on the {\it harmonic foliation} $\FF=\FF_S(A,B)$ connecting $A$ and $B$,
i.e.,  the gradient foliation of the harmonic function $\om$ vanishing on $0$ and equal to $1$ on $B$.
Hence $\WW_S(A,B)$ is equal to the $l_1$-norm of the associated  WAD $W_\FF$.

\subsection{Non-Intersection Principle}\label{so principle}

The following important principle says that
two wide quadrilaterals cannot go non-trivially one across the other:

\begin{lem}\label{no crossing}
Let us consider two quadrilaterals, $Q_1$ and $Q_2$,
endowed with the vertical foliations $\FF_1$ and $\FF_2$.
If $\WW(Q_i)\geq 1$ then \\
$\bullet$ either there exists a pair of disjoint leaves $\gamma_i$ of the foliations $\FF_i$;\\
$\bullet$ or $\WW(Q_1)=\WW(Q_2)=1$, and the rectangles perfectly match in the sense that
   the vertical sides of one of them coincide with the horizontal sides of the other.
\end{lem}


\begin{proof}
  Assume that the first option is violated,
so that, every leaf of $\FF_1$ crosses every leaf of $\FF_2$.
  Let us uniformize our rectangles by standard rectangles,
$\phi_i\from \BQ(a_i)\to Q_i$, where $a_i\geq 1$.
Let $\la_i$ be the Euclidean metrics on the $\BQ(a_i)$,
and let $\mu_i= (\phi_i)_*(\la_i)$.
Since every vertical leave $\gamma$ of $Q_2$ crosses every vertical leaf of $Q_1$,
$\mu_1(\gamma)\geq a_1$. Hence $a_2= \WW (Q_2)\leq 1/a_1$,
which implies that $a_1=a_2=1$. Thus, both $\BQ(a_i)\equiv \BQ$ are the squares.

Moreover, $\mu_1$ must be the extremal metric for $\FF_2$.
Since the extremal metric is unique (up to scaling), we conclude that
$\mu_1=\mu_2$. Hence $\phi_2^{-1}\circ \phi_1\from \BQ \to \BQ $
is the isometry of  the square, and the conclusion follows.
\end{proof}

